\chapter{Introduction}
\label{ch:introduction}

\section{Background}
\label{sec:background}

Large language models have moved from research curiosity to production infrastructure faster
than most enterprise architects anticipated. Within a few years of GPT-3's release, companies
across logistics, finance, healthcare, and manufacturing began deploying LLM-powered agents to
parse reports, answer supply chain queries, and generate operational summaries. The appeal is
obvious: a well-prompted model can synthesize hundreds of data points into a coherent narrative
in seconds, replacing workflows that once required a team of analysts several hours.

The economics, however, tell a more uncomfortable story. GPT-4o charges approximately
\$0.03 per thousand tokens, and enterprise deployments are not small. A company running a
thousand analytical queries per day, each fed a modest 5,000-token context window, spends
between \$450 and \$1,500 per month on token costs alone---before accounting for response
tokens, retries, or the larger contexts that complex queries require. At 10,000 queries per
day, the same arithmetic puts the bill at \$4,500 to \$15,000 monthly. These are not
hypothetical projections; they are the numbers that show up in cloud billing dashboards for
teams that have deployed LLM pipelines without a compression strategy.

The fundamental tension is architectural. LLMs are stateless: every call must receive the
full context needed to answer the query, because the model retains nothing between calls.
In a dynamic business environment, that context is the raw event stream---hundreds or
thousands of individual records describing sales, shipments, quality alerts, and contract
updates. Forwarding the full stream to every agent call is conceptually simple but
economically unsustainable.

This project addresses that tension by building a middleware layer that processes the raw
event stream before it reaches any LLM agent. The middleware uses a knowledge graph to
represent business state, computes structural deltas between successive snapshots, scores
each delta by its business significance, and assembles a compact Business Intelligence
Briefing that captures the most salient changes in a few hundred tokens. The result is a
system that lets LLM agents reason about enterprise data at a fraction of the token cost,
without meaningful loss of analytical quality.

\section{Motivation}
\label{sec:motivation}

Three converging developments make this project both necessary and achievable right now.

\textbf{LLM costs scale with context, not complexity.}
A language model does not charge more for a harder question; it charges more for a longer
prompt. This asymmetry means that a business analyst querying a 1,000-event data window
pays roughly ten times more than a colleague querying 100 events, even when the underlying
question is identical. Retrieval-Augmented Generation (RAG) addresses part of this by
fetching only relevant document chunks, but RAG assumes that the data can be indexed into
a vector store beforehand. Real-time event streams---where relevance is defined by what
changed in the last hour, not by semantic similarity to a query---do not fit neatly into
the RAG paradigm.

\textbf{Knowledge graphs have matured for structured business data.}
Libraries like NetworkX make it practical to build in-memory graph representations of
business entities and their relationships in a few dozen lines of Python. Graph-theoretic
measures such as degree centrality and betweenness centrality can quantify the relative
importance of each entity without any labeled training data. A supplier node that sits on
the shortest path between many product nodes has high betweenness centrality; when it
develops a delay, the effect propagates broadly and the node's centrality score provides
a natural measure of how much attention that delay deserves.

\textbf{Saliency-based filtering is proven in adjacent domains.}
Attention mechanisms in neural networks have demonstrated for years that not all input
tokens matter equally for a given output. In information retrieval, threshold-based
filtering of low-relevance documents routinely improves both latency and precision.
This project applies the same principle to structured business events: score each
change by its business impact, discard changes below a saliency threshold, and forward
only the high-impact changes to the agent. The contribution is to operationalize this
idea in the specific context of supply chain event streams and LLM agent pipelines.

\section{Project Scope}
\label{sec:scope}

This project is deliberately \textbf{software-focused}. No new hardware has been designed,
no proprietary dataset has been collected, and no claims are made about production readiness
in any specific enterprise environment. The system runs entirely on a standard laptop, using
a synthetic supply chain domain to provide a controllable and reproducible evaluation setting.

The scope covers four deliverables: a neuro-symbolic distillation pipeline implemented in
Python, a Streamlit web dashboard that exposes the pipeline interactively, a benchmark
framework that compares the pipeline against two baselines (naive full-context and
aggregation-only), and a ground truth evaluation against four injected anomaly patterns.
What falls outside the scope is equally important to state: no real enterprise data
connectors have been built, the saliency threshold was tuned by hand rather than learned
from labeled data, and the LLM integration currently relies on a single provider
(Google Gemini Flash) rather than a provider-agnostic abstraction.

The synthetic supply chain domain consists of ten event types spanning sales, restocks,
shipments, returns, price changes, supplier delays, quality issues, demand spikes, stockouts,
and new contracts. The knowledge graph uses a fixed topology of products, suppliers,
warehouses, and regions, giving the evaluation a well-defined ground truth against which to
measure the pipeline's detection performance. This controlled setup is a deliberate choice:
it allows precise measurement of compression ratios and recall without the confounding
variables that real enterprise data would introduce.

\section{Problem Statement}
\label{sec:problem}

The problem this project addresses is concrete: there is no open, benchmarked system that
combines symbolic preprocessing of business events with neural saliency scoring and
structured distillation into LLM-ready briefings. Four specific gaps define the current
state of the art.

\begin{enumerate}
  \item \textbf{No symbolic preprocessing layer for LLM agents.} Existing LLM deployments
  forward raw data---flat CSV rows, JSON event objects, or unstructured log lines---directly
  to the model prompt. The model is left to do all the relationship inference and change
  detection in-context, consuming tokens to reconstruct structure that could have been
  computed symbolically before the API call.

  \item \textbf{No saliency-aware context compression for business event streams.}
  RAG-based compression works well for document corpora but poorly for time-ordered event
  streams where relevance is determined by recency and structural change, not by semantic
  similarity to a query vector. No published system applies centrality-weighted saliency
  scoring to knowledge graph deltas for the purpose of LLM token reduction.

  \item \textbf{No supply-chain-specific multi-agent middleware.} Supply chain analysis
  requires reasoning about relationships---supplier-to-product, product-to-warehouse,
  warehouse-to-region---that a flat event list obscures. Existing agent frameworks provide
  tool-use and role-based routing, but they do not include a preprocessing layer that
  surfaces these relationships before the agent receives its context.

  \item \textbf{No open benchmarked system.} The few industrial papers that describe
  token-reduction pipelines for enterprise LLM use cases report results on proprietary
  datasets with no reproducible baseline. This project contributes an end-to-end open
  implementation with documented baselines and replicable benchmark scripts.
\end{enumerate}

\section{Objectives}
\label{sec:objectives}

To close the gaps identified above, this project pursues five concrete objectives:

\begin{enumerate}
  \item \textbf{Build a neuro-symbolic distillation pipeline} that processes raw supply chain
  events through a knowledge graph, a semantic delta engine, and a saliency scorer to produce
  structured Business Intelligence Briefings.

  \item \textbf{Achieve 10$\times$ or greater token compression} with at least 90\% recall
  relative to the naive full-context baseline, across event scales from 100 to 1,000 events.

  \item \textbf{Implement a Streamlit dashboard} with four pages---Pipeline Demo, Knowledge
  Graph, Benchmarks, and Agent Insights---giving non-technical users interactive access to
  the system.

  \item \textbf{Benchmark against a naive baseline and an aggregation-only baseline},
  measuring token usage, quality recall, pipeline latency, and cost savings at multiple
  event scales.

  \item \textbf{Validate detection quality on ground truth patterns}, injecting four known
  anomalies into the synthetic data and measuring whether the pipeline identifies them
  correctly, including cascade effects that aggregation-only approaches miss.
\end{enumerate}

Each subsequent chapter addresses one or more of these objectives directly, and
Chapter~\ref{ch:conclusion} revisits them to show that each has been met.

\section{Connection to AlcoWatch}
\label{sec:alcowatch}

This project has a sibling in the same codebase. AlcoWatch, a prior work by a colleague
at Amity University Tashkent, addresses drunk-driving prevention through wearable
physiological sensing and on-device inference on a Wear OS smartwatch
\citep{shaposhnikova2025alcowatch}. That system reads PPG, EDA, and skin temperature from
a wrist device, classifies driver stress in real time, and transmits the result over BLE
to an Arduino that adjusts the vehicle's cabin environment.

The conceptual connection to the present work is at the level of the middleware idea.
AlcoWatch distills a stream of raw physiological sensor samples into a single structured
estimate---the stress level---before passing it to the actuator. This project does the
same thing one abstraction layer up: it distills a stream of raw business events into a
structured briefing before passing it to an LLM agent. Both systems apply the principle
that downstream consumers should receive processed, salient signals rather than raw data.
The domain and the technology stack are completely different, but the architectural intuition
is shared. This connection is discussed further in Chapter~\ref{ch:conclusion} as part of
the broader claim that neuro-symbolic distillation is a pattern applicable across sensing
and data pipeline contexts.
